\structural{ЗАВДАННЯ}

Виконати дослідження для такої системи.

Потік літаків, що потребують посадки в аеропорті, пуассонівський з інтенсивністю 10~літаків за годину. В аеропорту є дві посадкові смуги. Літак, зробивши посадку на смугу, звільняє її через 35~хвилин. Якщо літак потребує посадки, коли усі смуги зайняті, він очікує на посадку. Якщо очікування на посадку триває більше ніж $70 \pm 10$~хвилин, то виникає необхідність дозаправки літака, на що аеропорту доведеться витратити $1000 \pm 200$~грошових одиниць. Після 140~хвилин очікування літак відправляється на посадку в інший аеропорт. За кожний літак, що здійснив посадку без очікування, аеропорт має прибуток 2000~гривень, а за кожний, що здійснив посадку з очікуванням, -- $1500 \pm 100$~гривень. Витрати на будівництво додаткової посадкової смуги складають 3\,000\,000~гривень. В аеропорту не може бути побудовано більше ніж 10~додаткових смуг.

Метою моделювання є визначення кількості посадкових смуг, при якій інтервал часу їх окупності буде мінімальним, якщо час окупності додаткових посадкових смуг складає
\[
T = \frac{S}{p},
\]
де $S$ -- витрати на побудову смуг, $p$ -- прибуток за один рік.

Для вищезазначеної системи:
\begin{enumerate}[label=\arabic*)]
\item розробити опис концептуальної моделі системи;
\item виконати формалізацію опису об'єкта моделювання в термінах стохастичної мережі Петрі;
\item розробити алгоритм імітації моделі дискретно-подійної системи у відповідності до побудованого формального опису;
\item для доведення коректності побудованого алгоритму виконати верифікацію алгоритму імітації;
\item визначити статистичні оцінки заданих характеристик моделі, що є метою моделювання;
\item провести експериментальне дослідження моделі;
\item інтерпретувати результати моделювання та сформулювати пропозиції щодо поліпшення функціонування системи;
\item зробити висновки щодо складності розробки моделі та алгоритму імітації на основі використаного формалізму, отриманих результатів моделювання та їх корисності.
\end{enumerate}
